\documentclass[11pt]{article}
\usepackage[margin=1in]{geometry}
\usepackage[T1]{fontenc}
\usepackage{amsmath,amssymb,amsthm,booktabs,array}
\usepackage[hidelinks]{hyperref}
\newtheorem{theorem}{Theorem}
\newtheorem{lemma}[theorem]{Lemma}
\newcommand{\lk}{\operatorname{Lk}}
\newcommand{\CAT}{\operatorname{CAT}}
\newcommand{\edge}{\mathbin{\text{---}}}
\title{A two-cube cubulation of the AIM free-by-cyclic example}
\author{Alec Kriebel\\Independent researcher\\\small ORCID: \href{https://orcid.org/0009-0001-9320-500X}{0009-0001-9320-500X}}
\date{October 12, 2026}
\begin{document}
\maketitle
\begin{abstract}
We construct a finite locally $\CAT(0)$ unit cube complex with one vertex,
seven edges, eight squares, and two three-dimensional cubes, whose fundamental
group is the mapping torus of $a\mapsto a$, $b\mapsto ab$, $c\mapsto bcb$.
Its universal cover gives a proper cocompact cubical action of the exact
group in AIM Question~5.2. All attaching maps, signed vertex links, and
reversible presentation changes are given explicitly. We also explain a
qualitative existence deduction from earlier WPD, Morse, and cubical
completion theorems. The contribution presented here is the explicit small
finite model; no historical firstness or minimality claim is made.
\end{abstract}

\section{The exact statement}

Let $F=F(a,b,c)$ and let $\phi\in\operatorname{Aut}(F)$ be given by
\[
 \phi(a)=a,\qquad \phi(b)=ab,\qquad \phi(c)=bcb.
\]
Indeed its inverse sends $a$ to $a$, $b$ to $a^{-1}b$, and $c$ to
$b^{-1}acb^{-1}a$; composition in either order freely reduces to the
respective basis generators. We use the convention
\begin{equation}\label{eq:G}
 G=F\rtimes_\phi\mathbb Z
 =\langle a,b,c,t\mid tat^{-1}=a,\ tbt^{-1}=ab,\ tct^{-1}=bcb\rangle.
\end{equation}
The final $b$ is positive. AIM Question~5.2, attributed to Rylee Lyman,
asks whether this group acts properly and cocompactly on a $\CAT(0)$ cube
complex~\cite{AIM}. The workshop report discusses the different formula
$b\mapsto aba$, $c\mapsto bcb$ and an earlier three-dimensional
candidate~\cite[p.~2]{workshop}. We keep that related example distinct.

\begin{theorem}\label{thm:main}
There is a finite locally $\CAT(0)$ unit cube complex $Z$ of dimension three
with cell counts $(1,7,8,2)$ and $\pi_1Z\cong G$. Consequently $G$ acts
freely, properly, and cocompactly by cubical isometries on the $\CAT(0)$
cube complex $\widetilde Z$.
\end{theorem}

Here a cube complex allows isometric identifications of boundary faces;
characteristic cubes in the quotient need not be embedded. This convention
is essential for the stated cell counts. The flag-link criterion for this
convention is Leary's Theorem~C.9~\cite[p.~33]{Leary}. His Theorem~C.4
(p.~31) ensures that a $\CAT(0)$ universal cover is cubical under the
stricter embedded-cube convention; Corollary~C.11 (p.~33) gives a strict
cubical quotient after the second cubical subdivision, if desired. The
counts in Theorem~\ref{thm:main} are before subdivision.

The construction uses a finite cyclic local isometry into a two-square
base. We verify the quotient links directly and identify the entire
fundamental group. No minimum dimension, minimum cell count, or virtual
specialness conclusion is asserted. Earlier tools also imply existence
through the target-specific deduction in Section~\ref{sec:older}. A bounded
source audit on October~12, 2026 found no prior documented exact answer or
directly covering geometric-cubulation theorem in the sources examined.
This is not a certificate of historical firstness or absolute novelty.

\section{The two-square base and cyclic domain}

In attaching words, capitals denote inverse oriented edges. Vertex names
$A,B,C,D$ are independent of this convention. Let $K$ have one vertex,
three oriented unit edges $x,y,b$, and two unit squares with boundary words
\begin{equation}\label{eq:K}
 xyXY,\qquad bxBY.
\end{equation}
Thus $H=\pi_1K=\langle x,y,b\mid [x,y]=1,\ bxb^{-1}=y\rangle$.
For an edge $h$, write $h^+$ and $h^-$ for the outgoing and incoming
germs. The link $L_K$ has six vertices and precisely the eight edges
\begin{equation}\label{eq:base-link}
 \{x^\epsilon,y^\delta\}\quad(\epsilon,\delta\in\{+,-\}),\qquad
 \{b^-,x^+\},\ \{b^-,x^-\},\ \{b^+,y^+\},\ \{b^+,y^-\}.
\end{equation}
It is a simple bipartite graph, with parts
$\{x^+,x^-,b^+\}$ and $\{y^+,y^-,b^-\}$. It has no triangles and is
therefore flag, so $K$ is locally $\CAT(0)$.

The assignments
\begin{equation}\label{eq:lambda}
 \lambda(x)=y,\qquad\lambda(y)=x,\qquad\lambda(b)=B
\end{equation}
extend to a cubical isometric involution of $K$. They take the two boundary
words in~\eqref{eq:K} to $yxYX$ and $BybX$, respectively. The first is a
dihedral reparametrization of the first square; the second is a cyclic
rotation of the inverse of the second boundary. The corresponding square
isometries agree with~\eqref{eq:lambda} on their faces. Their squares are
the identity, so they define the claimed involution, including on cells.

Define $Q$ by four distinct vertices $A,B,C,D$ and the following six unit
edges. The final column specifies a cubical map $j:Q\to K$.
\begin{center}
\begin{tabular}{cccc}
\toprule
edge & initial vertex & terminal vertex & $j$-label\\
\midrule
$p$ & $B$ & $A$ & $x$\\
$q$ & $A$ & $C$ & $y$\\
$r$ & $D$ & $C$ & $x$\\
$s$ & $B$ & $D$ & $y$\\
$d$ & $A$ & $B$ & $b$\\
$e$ & $C$ & $A$ & $b$\\
\bottomrule
\end{tabular}
\end{center}
Attach two squares with words
\begin{equation}\label{eq:Q}
 pqr^{-1}s^{-1}\quad\text{based at }B,
 \qquad dpe^{-1}q^{-1}\quad\text{based at }A.
\end{equation}
Their characteristic vertex cycles are $(B,A,C,D)$ and $(A,B,A,C)$.
In particular, the two occurrences of $A$ in the second square must both
be retained in corner calculations. The labels take these boundaries to
$xyXY$ and $bxBY$, so $j$ is defined on squares by unit-square isometries.

The complete signed links are given below. A path lists all link vertices
and all edges; there are no additional edges or higher simplices. The last
two columns record their images in $L_K$ under the two endpoint maps.
\begin{center}
\small
\setlength{\tabcolsep}{4pt}
\begin{tabular}{ccll}
\toprule
vertex & path in $\lk_Q$ & image under $j$ & image under $\lambda j$\\
\midrule
$A$ & $e^-\edge p^-\edge q^+\edge d^+$
 & $b^-\edge x^-\edge y^+\edge b^+$
 & $b^+\edge y^-\edge x^+\edge b^-$\\
$B$ & $d^-\edge p^+\edge s^+$
 & $b^-\edge x^+\edge y^+$
 & $b^+\edge y^+\edge x^+$\\
$C$ & $r^-\edge q^-\edge e^+$
 & $x^-\edge y^-\edge b^+$
 & $y^-\edge x^-\edge b^-$\\
$D$ & $r^+\edge s^-$
 & $x^+\edge y^-$
 & $y^+\edge x^-$\\
\bottomrule
\end{tabular}
\end{center}
For example, the two $A$-corners of the second square supply the distinct
edges $\{d^+,q^+\}$ and $\{p^-,e^-\}$, while the first square supplies
$\{p^-,q^+\}$. Every row maps injectively to $L_K$. Comparison
with~\eqref{eq:base-link} shows that each image is a full induced
subcomplex: every base-link edge between listed vertices is listed.
Consequently $Q$ is locally $\CAT(0)$, and $j$ and $\lambda j$ are
cubical local isometries, by the injective-full-link characterization.

The edges $p,q,s$ form a spanning tree. On collapsing them, the first
square kills $r$ and the second gives $d=e$. Hence
\begin{equation}\label{eq:cyclic}
 \pi_1(Q,A)=\langle dp\rangle\cong\mathbb Z,
 \qquad j_*(dp)=bx,
 \qquad (\lambda j)_*(dp)=By.
\end{equation}
The maps are injective on fundamental groups; alternatively $bx$ has
infinite order because the homomorphism $H\to\mathbb Z$ sending $b$ to
$1$ and $x,y$ to $0$ sends $bx$ to $1$.

\section{The finite quotient and its flag link}

Form the double mapping cylinder
\begin{equation}\label{eq:Z}
 Z=K\cup_{j,\lambda j}\bigl(Q\times[0,1]\bigr),
\end{equation}
where $(q,0)$ is attached by $j$ and $(q,1)$ by $\lambda j$. Give the
interval length one and the product its cubical structure. Face
identifications come from the genuine product and the two cubical maps.
They need not be globally injective. Denote the vertical edges by
$z_A,z_B,z_C,z_D$, directed from the $j$ end to the $\lambda j$ end. These
are four distinct edges, although all their endpoints become the sole
vertex of $K$. The second product cube has $z_A$ at two distinct boundary
positions, just as the second square of $Q$ has two $A$-corners.

There are three base edges and four vertical edges, two base squares and
six edge-product squares, and two square-product three-cubes. Thus
\begin{equation}\label{eq:counts}
 (\#Z^{(0)},\#\text{edges},\#\text{squares},\#\text{three-cubes})=(1,7,8,2).
\end{equation}
The Euler characteristic is $1-7+8-2=0$.

Let $L$ be the unique vertex link of $Z$. It consists of $L_K$ and eight
cones with apices $z_v^+$ and $z_v^-$, $v\in\{A,B,C,D\}$. Their bases
are the $j$ and $\lambda j$ link images in the table, respectively. This
description follows corner by corner from the product: an edge-product
square adds an apex--base edge, and a square-product cube adds the cone
triangle over each corresponding square corner. No two apices are
adjacent, since a product-cube corner has only one interval direction.
The eight apices are distinct germs of the four vertical edges.

The link maps in the table are injective, so no cone has repeated vertices,
edges, or triangles. Distinct cones have distinct apices. At a repeated
$A$-corner the two triangles from the second product cube have distinct
base edges, as displayed above; they are not a repeated simplex. Therefore
$L$ is an actual simplicial complex, not merely a graph with prescribed
clique counts.

If a clique of $L$ contains no apex, it lies in the flag complex $L_K$.
Otherwise it contains exactly one apex. Its other vertices lie in the
full cone base, and their base-link clique spans a simplex in that base.
Adjoining the apex gives a simplex in the cone. Thus $L$ is flag. Its
counts provide an additional check: it has $6+8=14$ vertices,
$8+2(4+3+3+2)=32$ edges, and $2(3+2+2+1)=16$ triangles, with no
four-clique. Leary's general cube-complex criterion~\cite[Theorem~C.9,
p.~33]{Leary} now shows that $Z$ is locally $\CAT(0)$ and that its
universal cover is $\CAT(0)$.

\section{The fundamental group and geometric action}

We record a direct presentation computation to verify the exact group,
including the sign of the last $b$. Besides~\eqref{eq:K}, the six
edge-product squares have relators
\begin{equation}\label{eq:squares}
\begin{aligned}
 &xz_A Yz_B^{-1}, &&yz_C Xz_A^{-1}, &&xz_C Yz_D^{-1},\\
 &yz_D Xz_B^{-1}, &&bz_Bbz_A^{-1}, &&bz_Abz_C^{-1}.
\end{aligned}
\end{equation}
Attaching three-cells does not change the fundamental group. Write
$z=z_A$. The $p,q,s$ relations eliminate the other vertical generators:
\begin{equation}\label{eq:eliminate}
 z_B=xzY,\qquad z_C=Yzx,\qquad z_D=YxzYx.
\end{equation}
The $r$ relation is then redundant using $xy=yx$. The $d$ relation is
$z^{-1}(bx)z=By$; the $e$ relation is $z^{-1}(yb)z=xB$. These agree,
because $bxB=y$ gives $yb=bx$ and $xB=By$. Each elimination is a
reversible Tietze transformation, so
\begin{equation}\label{eq:HNN}
 \pi_1Z=\langle x,y,b,z\mid xy=yx,\ bxB=y,\ z^{-1}(bx)z=By\rangle.
\end{equation}

Make the invertible change of generators
\begin{equation}\label{eq:basis}
 t=x,\quad a=xy^{-1},\quad c=(zx)^{-1};
 \qquad x=t,\quad y=a^{-1}t,\quad z=c^{-1}t^{-1},
\end{equation}
leaving $b$ fixed. The first two relations of~\eqref{eq:HNN} become
$[a,t]=1$ and $btb^{-1}=a^{-1}t$. The latter is equivalent to
$tbt^{-1}=ab$. For the third relation put $u=Bx$, $v=bx$. Since
$By=xB=xuX$, its successive equivalent forms are
\begin{align*}
 z^{-1}vz=xux^{-1}
 &\quad\Longleftrightarrow\quad (zx)u(zx)^{-1}=v\\
 &\quad\Longleftrightarrow\quad c^{-1}(b^{-1}t)c=bt\\
 &\quad\Longleftrightarrow\quad tc=bcbt\\
 &\quad\Longleftrightarrow\quad tct^{-1}=bcb.
\end{align*}
Thus~\eqref{eq:HNN} becomes exactly~\eqref{eq:G}. Explicit inverse
substitutions and reversible relations establish an isomorphism, rather
than only a homomorphism obtained by checking relators in $G$.

The cover $\widetilde Z$ is three-dimensional with finite links and unit
Euclidean cube shapes. It is locally finite; its natural complete length
metric is proper (equivalently, local compactness and completeness give
properness for this geodesic metric). The deck group acts freely by
cubical isometries. Covering neighborhoods show that this action is
properly discontinuous, and, on the proper cover, that it is a proper
action: only finitely many translates of a compact set meet that set.
Its quotient is the finite complex $Z$, so the action is cocompact.
The isomorphism above identifies this deck group with the full group $G$.
This proves Theorem~\ref{thm:main}.

\section{An existence deduction from earlier tools}\label{sec:older}

We give a second route to qualitative existence to clarify the relation
to prior methods. This route obtains an unspecified finite cyclic local
isometry by existing completion theorems; it does not supply the cell
counts~\eqref{eq:counts}.

Write $A_0=\langle x,y\mid[x,y]=1\rangle\cong\mathbb Z^2$ and
$H=\langle A_0,b\mid bxb^{-1}=y\rangle$. The HNN normal-form theorem
embeds $A_0$. In its Bass--Serre tree, vertices are left cosets $hA_0$,
and an oriented edge $e$ goes from $A_0$ to $bA_0$, with stabilizer
$\langle y\rangle$. Set $g=bx$. The homomorphism
$\epsilon_b:H\to\mathbb Z$, with $\epsilon_b(b)=1$ and
$\epsilon_b(x)=\epsilon_b(y)=0$, gives a height on vertices. Since
$gA_0=bA_0$, the edges $g^ne$ concatenate, increasing height at every
step. They form a nonbacktracking bi-infinite axis of $g$.
At $A_0$, the outgoing edge has stabilizer $\langle y\rangle$ and the
incoming edge has stabilizer
\[
 g^{-1}\langle y\rangle g
 =x^{-1}b^{-1}\langle y\rangle bx=\langle x\rangle.
\]
Their intersection is trivial in $A_0$. Hence the pointwise stabilizer
of the axis segment from $g^{-1}A_0$ to $gA_0$ is trivial.
Minasyan--Osin~\cite[Corollary~4.3]{MO}, in the corrected v6, gives WPD
for this specific loxodromic $g$. No finite-extension claim is used.

Dahmani--Guirardel--Osin~\cite[Theorem~6.8 and Corollary~2.9]{DGO}
provide a virtually cyclic hyperbolically embedded subgroup $E(g)$
containing $g$. As $\epsilon_b(g)=1$, $g$ has infinite order. Sisto's
Theorem~1~\cite{Sisto} then makes $g$ Morse; in particular
$C=\langle g\rangle$ is quasiconvex in the word metric on $H$.
The metric required for completion matches exactly: $K$ has one vertex,
so $\widetilde K^{(1)}$ is the Cayley graph for $\{x,y,b\}$, with its
edge metric. The two-edge circle $\Gamma\to K$ reading $bx$ realizes
the inclusion $C\hookrightarrow H$, and its lifted image contains $C$
and lies within distance one of it. Geodesics between its vertices also
stay uniformly close, by adding endpoint paths of length at most one
and using the Morse property. Thus no conversion from CAT(0)-metric
quasiconvexity is being assumed.

Ben-Zvi--Kropholler--Lyman~\cite[Theorems~1.1 and 5.1]{BKL} now give a
finite pointed locally $\CAT(0)$ cube complex $Q_0$ and a cubical local
isometry $j_0:Q_0\to K$, through which $\Gamma\to K$ factors, whose
fundamental-group image is exactly $C$. Their maximal-tree generating
set is $\{x,y,b\}$ here. Local-isometry injectivity gives
$\pi_1Q_0\cong C$, retaining the generator $bx$, rather than a power of
it. The original $bx$ circle itself is not a local isometry: at its
intermediate vertex the germs $b^-,x^+$ span a base-link edge. Completion
is what repairs that defect.

Postcomposition with the explicit involution~\eqref{eq:lambda} gives
the second local isometry $\lambda j_0$ on the \emph{same} cubical
domain, taking the generator to $By$. Hence
$M=K\cup_{j_0,\lambda j_0}(Q_0\times[0,1])$ is finite. Its links
have the same general structure as those of $Z$: each interval germ
cones off an injective full link image, and distinct interval germs are
not adjacent. The clique argument in Section~3 proves that these links
are flag. This also gives a direct verification of the local-isometry
gluing discussed in Hagen--Wise~\cite[Remark~3.7]{HW}. Because a local
isometry into $K$ has domain dimension at most two, $\dim M\leq3$.
Van Kampen gives~\eqref{eq:HNN}, and~\eqref{eq:basis} identifies
$\pi_1M$ with $G$. Its finite locally $\CAT(0)$ quotient therefore
gives another proper cocompact action of the exact group.

The abstract tools in this deduction are prior work. Their availability
does not date the discovery or publication of this exact application.
The present note makes the element, metric, compatible maps, and group
identification explicit. Hagen--Wise's current general fast-automorphism
result~\cite[Corollary~7.2]{HW} supplies free cubical actions; properness
and cocompactness require the additional argument above or the explicit
model of Sections~2--4.

\section*{Verification and disclosure}

The accompanying verification package contains three Python standard-library
programs for finite combinatorial and algebraic checks. They check the
specified face maps, signed links and fullness, the involution, the
fourteen-vertex link, cell counts, and generator substitutions, with
validation active in ordinary and optimized Python execution. These
finite checks supplement the written proof; they do not constitute a
machine proof of the global CAT(0) or deck-action conclusions.

Extensive AI assistance was used in discovery, drafting, and internal
independent adversarial checks. This note is unrefereed; there has been
no external human peer review and no proof-assistant certification. The
author is responsible for the stated claims and their scope.

\begin{thebibliography}{9}
\small
\bibitem{AIM}
AIM Problem Lists, \emph{Rigidity properties of free-by-cyclic groups},
Section~5, Question~5.2 (Rylee Lyman),
\url{http://www.aimpl.org/freebycyclic/5/}.

\bibitem{workshop}
American Institute of Mathematics, \emph{Rigidity properties of
free-by-cyclic groups: workshop summary}, workshop October~2--6, 2023,
\url{https://aimath.org/pastworkshops/freebycyclicrep.pdf}, p.~2.

\bibitem{BKL}
M.~Ben-Zvi, R.~Kropholler, and R.~A.~Lyman,
\emph{Folding-like techniques for CAT(0) cube complexes},
Math. Proc. Cambridge Philos. Soc. \textbf{173} (2022), 227--238,
\href{https://doi.org/10.1017/S0305004121000645}{doi:10.1017/S0305004121000645};
\href{https://arxiv.org/abs/2011.05374v1}{arXiv:2011.05374v1},
Theorems~1.1 and 5.1.

\bibitem{DGO}
F.~Dahmani, V.~Guirardel, and D.~Osin,
\emph{Hyperbolically embedded subgroups and rotating families in groups
acting on hyperbolic spaces}, Mem. Amer. Math. Soc. \textbf{245} (2017),
no.~1156; \href{https://arxiv.org/abs/1111.7048v5}{arXiv:1111.7048v5},
Theorem~6.8 and Corollary~2.9.

\bibitem{HW}
M.~Hagen and D.~T.~Wise,
\emph{Cubulating mapping tori of some polynomial growth free group
automorphisms}, \href{https://arxiv.org/abs/1605.07879v3}{arXiv:1605.07879v3}
(2025), Remark~3.7 and Corollary~7.2.

\bibitem{Leary}
I.~J.~Leary, \emph{A metric Kan--Thurston theorem},
\href{https://arxiv.org/abs/1009.1540v3}{arXiv:1009.1540v3} (2012),
Appendix~C: Theorem~C.4, p.~31; Theorem~C.9 and Corollary~C.11, p.~33.

\bibitem{MO}
A.~Minasyan and D.~Osin,
\emph{Acylindrical hyperbolicity of groups acting on trees},
\href{https://arxiv.org/abs/1310.6289v6}{arXiv:1310.6289v6} (2017,
corrected version), Corollary~4.3.

\bibitem{Sisto}
A.~Sisto, \emph{Quasi-convexity of hyperbolically embedded subgroups},
\href{https://arxiv.org/abs/1310.7753v1}{arXiv:1310.7753v1} (2013),
Theorem~1.
\end{thebibliography}
\end{document}
